\begin{abstract}
SO$(N)$ Chern--Simons 理論の彩色 Kauffman 不変量の large $N$ 双対（\cite{Marino09}）を、\cite{AgSc12} の orientifold の refined 位相的弦の枠組みで refined にし、refined クロスキャップ不変量を定義してその整数性・構造・正値性を調べる。
定義は、SO$(2N)$ の refined Chern--Simons 不変量を双対 Cauchy 核（t-brane）で生成関数にまとめ、被覆 U$(2N-1)$ の composite 表現の寄与の半分を引く、というものである：$g = f - \tfrac12 f^{\cov}$。規格化は、行列式表現と荷電共役で不変になる一様な規格化をとる。

主な結果は次のとおりである。
\begin{enumerate}
\item \textbf{整数性と普遍因子}：DAHA-Jones 多項式を一般の $t$ で計算し、trefoil の 2 箱・3 箱と $T(2,5)$ の 2 箱の refined クロスキャップ不変量を一般の $N$ で求めた。これらは整数係数の Laurent 多項式で、$N$ に依存する極はすべて消え、$t = q$ で Mariño の結果に戻る。refined だけの部分は、表現 [1]・2 箱・3 箱・Hopf 絡み目のすべてで普遍因子 $(\bt^2 - 1)(\ba^2 - \bq^2)$ で割り切れる。
\item \textbf{全 $d$ の構造}：orientifold の M 理論の指数で一粒子状態を $\tilde I$ で入れ替わる対・不変な状態・凍結した状態に分けると、$g = \tfrac12 E/(1+t) + C$ が従う。$E$ は brane と image の箱が同数のセクターにしか現れず、$(1+t)$ の極の位置がこの予言と一致する。境界ゼロモードへの $\tilde I$ の作用から $E$ の境界因子 $z_q^{n-1}(z_q^+)^n$ を導き、Hopf 絡み目・結び目・framed unknot で確かめた。
\item \textbf{表現 [1] の正値性}：トーラス結び目 $T(2,2m+1)$（$m \le 5$）の表現 [1] では、$c = 1$ 部分が $G$ 展開でき、$\lambda^\beta$ ごとに符号がそろい、$c = 2$ 部分は $-(\ba^2 - \bq^2)(\bt^2 - 1)\times$ 正の多項式になる。これを判定式として予想の形にまとめる。
\item \textbf{被覆の構造}：被覆を brane と image に分けると、brane--brane のセクターは \cite{KN17} の U$(N_A)$ の refined LMOV 不変量に全係数で一致する。brane--image のセクターは、境界因子を外すと $\mathrm{Sym}^k V$（$V$ は trefoil の reduced superpolynomial）、$S \otimes S$、普遍的な補正から成る状態の空間として同定でき、14 のトーラス結び目で行の母関数が成り立つ。この同定から $\tilde I$ のトレース $E$ が条件なしに作れる。
\item \textbf{refined Sp$(2N)$}：C$_N$ 型の DAHA を実装し、\cite{AM22} の refined Sp$(2N)$ の結び目不変量を計算した。表現 [1] では $g^{\mathrm{Sp}}(\ba) = -g^{\SO}(-\ba)$ が、2 箱では $W^{\mathrm{Sp}}_R(\ba) = (-1)^{|R|}W^{\SO}_R(-\ba)$ が成り立つ。
\end{enumerate}
多箱のクロスキャップ部分では、$c = 1$ 部分が $\sigma$ 対称でなく $G$ 展開できず、符号もそろわない。被覆側に問題が無いことから、障害は凍結した状態の多重項の構造（向き付けられない曲面の refined のハンドルの因子）に絞られる。これを未解決問題として提示する。
\end{abstract}
